\documentclass[final]{beamer}\usepackage[]{graphicx}\usepackage[]{color}
%% maxwidth is the original width if it is less than linewidth
%% otherwise use linewidth (to make sure the graphics do not exceed the margin)
\makeatletter
\def\maxwidth{ %
  \ifdim\Gin@nat@width>\linewidth
    \linewidth
  \else
    \Gin@nat@width
  \fi
}
\makeatother

\definecolor{fgcolor}{rgb}{0.345, 0.345, 0.345}
\newcommand{\hlnum}[1]{\textcolor[rgb]{0.686,0.059,0.569}{#1}}%
\newcommand{\hlstr}[1]{\textcolor[rgb]{0.192,0.494,0.8}{#1}}%
\newcommand{\hlcom}[1]{\textcolor[rgb]{0.678,0.584,0.686}{\textit{#1}}}%
\newcommand{\hlopt}[1]{\textcolor[rgb]{0,0,0}{#1}}%
\newcommand{\hlstd}[1]{\textcolor[rgb]{0.345,0.345,0.345}{#1}}%
\newcommand{\hlkwa}[1]{\textcolor[rgb]{0.161,0.373,0.58}{\textbf{#1}}}%
\newcommand{\hlkwb}[1]{\textcolor[rgb]{0.69,0.353,0.396}{#1}}%
\newcommand{\hlkwc}[1]{\textcolor[rgb]{0.333,0.667,0.333}{#1}}%
\newcommand{\hlkwd}[1]{\textcolor[rgb]{0.737,0.353,0.396}{\textbf{#1}}}%

\usepackage{framed}
\makeatletter
\newenvironment{kframe}{%
 \def\at@end@of@kframe{}%
 \ifinner\ifhmode%
  \def\at@end@of@kframe{\end{minipage}}%
  \begin{minipage}{\columnwidth}%
 \fi\fi%
 \def\FrameCommand##1{\hskip\@totalleftmargin \hskip-\fboxsep
 \colorbox{shadecolor}{##1}\hskip-\fboxsep
     % There is no \\@totalrightmargin, so:
     \hskip-\linewidth \hskip-\@totalleftmargin \hskip\columnwidth}%
 \MakeFramed {\advance\hsize-\width
   \@totalleftmargin\z@ \linewidth\hsize
   \@setminipage}}%
 {\par\unskip\endMakeFramed%
 \at@end@of@kframe}
\makeatother

\definecolor{shadecolor}{rgb}{.97, .97, .97}
\definecolor{messagecolor}{rgb}{0, 0, 0}
\definecolor{warningcolor}{rgb}{1, 0, 1}
\definecolor{errorcolor}{rgb}{1, 0, 0}
\newenvironment{knitrout}{}{} % an empty environment to be redefined in TeX

\usepackage{alltt}
\usepackage{grffile}
\mode<presentation>{\usetheme{CambridgeUSPOL}}

\usepackage[utf8]{inputenc}
\usepackage{amsfonts}
\usepackage{amsmath}
\usepackage{natbib}
\usepackage{graphicx}
\usepackage{array,booktabs,tabularx}
\newcolumntype{Z}{>{\centering\arraybackslash}X}

% rysunki
\usepackage{tikz}
\usepackage{ifthen}
\usepackage{xxcolor}
\usetikzlibrary{arrows}
\usetikzlibrary[topaths]
\usetikzlibrary{decorations.pathreplacing}
%\usepackage{times}\usefonttheme{professionalfonts}  % times is obsolete
\usefonttheme[onlymath]{serif}
\boldmath
\usepackage[orientation=portrait,size=a0,scale=1.4,debug]{beamerposter}                       % e.g. for DIN-A0 poster
%\usepackage[orientation=portrait,size=a1,scale=1.4,grid,debug]{beamerposter}                  % e.g. for DIN-A1 poster, with optional grid and debug output
%\usepackage[size=custom,width=200,height=120,scale=2,debug]{beamerposter}                     % e.g. for custom size poster
%\usepackage[orientation=portrait,size=a0,scale=1.0,printer=rwth-glossy-uv.df]{beamerposter}   % e.g. for DIN-A0 poster with rwth-glossy-uv printer check
% ...
%

\usecolortheme{seagull}
\useinnertheme{rectangles}
\setbeamercolor{item projected}{bg=darkred}
% \setbeamertemplate{enumerate items}[default]
\setbeamertemplate{caption}{\insertcaption} 
\setbeamertemplate{navigation symbols}{}
\setbeamercovered{transparent}
\setbeamercolor{block title}{fg=darkred}
\setbeamercolor{local structure}{fg=darkred}

\setbeamercolor*{enumerate item}{fg=darkred}
\setbeamercolor*{enumerate subitem}{fg=darkred}
\setbeamercolor*{enumerate subsubitem}{fg=darkred}

\setbeamercolor*{itemize item}{fg=darkred}
\setbeamercolor*{itemize subitem}{fg=darkred}
\setbeamercolor*{itemize subsubitem}{fg=darkred}

\newlength{\columnheight}
\setlength{\columnheight}{90cm}
\renewcommand{\thetable}{}
\def\andname{,}
\authornote{}

\renewcommand{\APACrefatitle}[2]{}
\renewcommand{\bibliographytypesize}{\footnotesize} 
\renewcommand{\APACrefYearMonthDay}[3]{%
  {\BBOP}{#1}
  {\BBCP}
}
\IfFileExists{upquote.sty}{\usepackage{upquote}}{}
\begin{document}






\date{}
\author{Piotr Sobczyk\inst{1}*, Micha\l{}  Burdukiewicz\inst{2}, Chris Lauber\inst{3}, Pawe\l{} Mackiewicz\inst{2} \\
*Piotr.Sobczyk@pwr.edu.pl}


\institute{

\textsuperscript{1}Wroc\l{}aw University of Technology, Institute of Mathematics and Computer Science, Poland

\vspace{0.3cm}

\textsuperscript{2}University of Wroc\l{}aw, Department of Genomics, Poland 

\vspace{0.3cm}

\textsuperscript{3}Dresden University of Technology, Institute of Medical Informatics and Biometry, Poland 

}
\title{\huge Quick Permutation Test: feature filtering of n-gram data}

\begin{frame}
  \begin{columns}
    \begin{column}{.49\textwidth}
      \begin{beamercolorbox}[center,wd=\textwidth]{postercolumn}
        \begin{minipage}[T]{.95\textwidth}
          \parbox[t][\columnheight]{\textwidth}
            {
    
        
    \begin{block}{Introduction}
      N-grams (k-tuples) are vectors of n characters derived from input sequence(s). They may form continuous sub-sequences or be discontinuous. Another important n-gram parameter is its position. Instead of just counting n-grams, one may want to count how many n-grams occur at a given position in multiple (e.g. related) sequences.

\small{
       \begin{columns}[c] % the "c" option specifies center vertical alignment
    \column{.5\textwidth} 
% latex table generated in R 3.1.3 by xtable 1.7-4 package
% Tue Apr  7 12:30:44 2015
\begin{table}[ht]
\centering
\begin{tabular}{rrrrrrr}
  \hline
 & P1 & P2 & P3 & P4 & P5 & P6 \\ 
  \hline
S1 & 4 & 3 & 1 & 3 & 1 & 3 \\ 
  S2 & 2 & 3 & 1 & 1 & 4 & 4 \\ 
  S3 & 4 & 3 & 1 & 3 & 4 & 4 \\ 
   \hline
\end{tabular}
\caption{Sample sequences.} 
\end{table}

      
      
     % column designated by a command

    \column{.5\textwidth}
    
% latex table generated in R 3.1.3 by xtable 1.7-4 package
% Tue Apr  7 12:30:44 2015
\begin{table}[ht]
\centering
\begin{tabular}{rrrr}
  \hline
1 & 2 & 3 & 4 \\ 
  \hline
2 & 0 & 3 & 1 \\ 
  2 & 1 & 1 & 2 \\ 
  1 & 0 & 2 & 3 \\ 
   \hline
\end{tabular}
\caption{Unigram counts.} 
\end{table}


    \end{columns}

% latex table generated in R 3.1.3 by xtable 1.7-4 package
% Tue Apr  7 12:30:44 2015
\begin{table}[ht]
\centering
\begin{tabular}{rrrrrrrrrrr}
  \hline
P1\_1 & P2\_1 & P3\_1 & P4\_1 & P5\_1 & P6\_1 & P1\_2 & P2\_2 & P3\_2 & P4\_2 & P5\_2 \\ 
  \hline
0 & 0 & 1 & 0 & 1 & 0 & 0 & 0 & 0 & 0 & 0 \\ 
  0 & 0 & 1 & 1 & 0 & 0 & 1 & 0 & 0 & 0 & 0 \\ 
  0 & 0 & 1 & 0 & 0 & 0 & 0 & 0 & 0 & 0 & 0 \\ 
   \hline
\end{tabular}
\caption{A fraction of possible unigrams with position information.} 
\end{table}

}      
    \end{block}
    
    \vfill
    
    \begin{block}{Curse of dimensionality}
Number of possible positioned n-grams (not taking into account distances between elemnts of n-gram):
$$
n_{max} = L \times m^n
$$

\centering
\scalebox{0.85}{  
\begin{knitrout}
\definecolor{shadecolor}{rgb}{0.969, 0.969, 0.969}\color{fgcolor}

{\centering \includegraphics[width=\maxwidth]{figure/unnamed-chunk-4-1} 

}



\end{knitrout}
}


    \end{block}
    \vfill
    
    
    \begin{block}{Permutation test}
      During permutation tests class labels are randomly exchanged during computation of significance statistic. p-values are defined as:
$$\textnormal{p-value} = \frac{N_{T_P > T_R}}{N}$$
where $N_{T_P > T_R}$ is number of times when $T_P$ (permuted test statistic) was more extreme than $T_R$ (test statistic for non-permuted data).

Permutation tests are model and statistic independent, but computationally expensive (especially precise estimation of low p-values, because the number of permutations is inversely proportional to the interval between p-values).
      
    \end{block}
    \vfill
    
    \begin{block}{QuiPT algorithm}
      If probability that target equals 1 is $p$ and probability that feature equals
1 is $q$ and feature and target are independent then each of them has the 
following probabilities 
$$P(Target, Feature) = (1,1)) = p \cdot q$$
$$P(Target, Feature) = (1,0)) = p \cdot (1-q)$$
$$P(Target, Feature) = (0,1)) = (1-p) \cdot q$$
$$P(Target, Feature) = (0,0)) = (1-p) \cdot (1-q)$$

\begin{equation*}
\begin{split}
F(n_{1,1}, n_{1,0}, n_{0,1}, n_{0,0}) = {n \choose n_{1,1}} (p\cdot q)^{n_{1,1}} {n - n_{1,1} \\ \choose n_{1,0}} (p\cdot (1-q))^{n_{1,0}} \\
{n - n_{1,1} - n_{1,0} \\ \choose n_{0,1}} ((1-p)\cdot q)^{n_{0,1}} \\
{n - n_{1,1} - n_{1,0} -n_{0,1}\choose n_{0,0}} \\((1-p)\cdot (1-q))^{n_{0,0}}$$
\end{split}
\end{equation}

In addition to this: 
$n_{1,\cdot} = n_{1,1} + n_{1,0}$ and
$n_{\cdot, 1} = n_{1,1} + n_{0,1}$ are known and fixed.
    \end{block}
    \vfill
    

    
            }
        \end{minipage}
      \end{beamercolorbox}
    \end{column}
    
    
%new column ------------------------------------------------------    
    
    \begin{column}{.52\textwidth}
      \begin{beamercolorbox}[center,wd=\textwidth]{postercolumn}
        \begin{minipage}[T]{.95\textwidth}  
          \parbox[t][\columnheight]{\textwidth}
            {
            
    
    \begin{block}{Validation procedure}
    \begin{enumerate}[1.]
      \item Chose randomly (without replacement) 3816 proteins without signal peptides, reshuffle 3816 proteins with signal peptides.
      \item Perform 5-fold cross-validation.
      \item Repeat step 1. and 2. 250 times.
    \end{enumerate} 
    \end{block}
    \vfill
    
    \begin{block}{Validation results}
    \begin{figure}
    \scalebox{0.85}{    
\begin{knitrout}
\definecolor{shadecolor}{rgb}{0.969, 0.969, 0.969}\color{fgcolor}\begin{kframe}


{\ttfamily\noindent\bfseries\color{errorcolor}{\#\# Error in t(poster\_data[["{}metrics"{}]][c("{}AUC"{}, "{}H"{}, "{}Gini"{}, "{}Recall"{}, "{}Spec"{}, : object 'poster\_data' not found}}

{\ttfamily\noindent\bfseries\color{errorcolor}{\#\# Error in ncol(metrics): object 'metrics' not found}}

{\ttfamily\noindent\bfseries\color{errorcolor}{\#\# Error in is.factor(x): object 'melted\_metrics' not found}}

{\ttfamily\noindent\bfseries\color{errorcolor}{\#\# Error in levels(melted\_metrics[["{}metric"{}]]) <- c("{}AUC"{}, "{}H-measure"{}, "{}Gini"{}, : object 'melted\_metrics' not found}}

{\ttfamily\noindent\bfseries\color{errorcolor}{\#\# Error in ggplot(melted\_metrics, aes(x = metric, y = value)): object 'melted\_metrics' not found}}\end{kframe}
\end{knitrout}
    }
    \end{figure}

The mean AUC yielded by cross-validation is  
    
    \begin{figure}
    \scalebox{0.85}{    
\begin{knitrout}
\definecolor{shadecolor}{rgb}{0.969, 0.969, 0.969}\color{fgcolor}\begin{kframe}


{\ttfamily\noindent\bfseries\color{errorcolor}{\#\# Error in eval(expr, envir, enclos): object 'poster\_data' not found}}

{\ttfamily\noindent\bfseries\color{errorcolor}{\#\# Error in eval(expr, envir, enclos): object 'position\_data' not found}}

{\ttfamily\noindent\bfseries\color{errorcolor}{\#\# Error in ggplot(position\_data, aes(x = Var1, y = Freq)): object 'position\_data' not found}}\end{kframe}
\end{knitrout}
    }
    \end{figure}
    \end{block}
    \vfill
    
    \begin{block}{Comparision with other signal peptide predictors}
    Benchmark data set: 140 eukaryotic proteins with signal peptide and 280 randomly chosen eukaryotic proteins without signal peptide added after 2010.
    
    \vspace{.5cm}
    
    signal.hsmm1987: trained on data set of 496 eukaryotic proteins with signal peptides added before year 1987.
    
    \vspace{.5cm}
    
    signal.hsmm2010: trained on data set of 3676 eukaryotic proteins with signal peptides added before year 2010.
    \end{block}
    \vfill
    
    
    
    \begin{block}{Comparision of various software}
    STH
    \end{block}
    \vfill
     
    \begin{block}{Summary}
      Hidden semi-Markov models can be used to accurately predict the presence of secretory signal peptides effectively extracting information from very small data sets.
    \end{block}
    \vfill 
    
        \begin{block}{Avaibility}
        \footnotesize{
      signal.hsmm web server: 
      
      \url{http://michbur.shinyapps.io/signalhsmm/}$\begin{array}{l}
        \begin{figure}[ht]
        \centering
        \scalebox{0.3}{
          \includegraphics{qrsignalhsmm.png}
        }
        \end{figure}
      \end{array}
      $
      
      signal.hsmm R package: 
      
      \url{http://cran.r-project.org/web/packages/signal.hsmm/}
        }
    \end{block}
    \vfill 
     
     
    \begin{block}{Bibliography}
    \tiny{
      \bibliographystyle{apalike}
      \bibliography{signalhsmm_poster}
    }
    \end{block}
    \vfill
            }
        \end{minipage}
      \end{beamercolorbox}
    \end{column}
  \end{columns}  
\end{frame}
\end{document}
